\documentclass{article}
\usepackage{graphicx}

\usepackage[utf8]{inputenc}     % pdfLaTeX
\usepackage[T1]{fontenc}
\usepackage[magyar]{babel}

\title{Az Ember}
\author{Haydu Lőrinc, Lengyel Zoltán Péter, Rovenszky Robin}
\date{Legutóbb frissítve: 2026. Április 22.}

\begin{document}

\maketitle
\tableofcontents
\newpage

\section{Az ember testfelépítése}
\section{A vér} % van PPT
\section{Az immunrendszer} % van PPT
\section{A keringési rendszer}
\section{A légzőrendszer}
\section{A táplálkozás}
\section{Az egészséges táplálkozás}
\section{A kiválasztórendszer}
\section{A mozgás}
\section{A kültakaró}
\section{Az életfolyamatok szabályozása} % maybe rename to Szabályozás
\section{A hormonális szabályozás}
% I and II
\section{Az idegrendszer}
% működése majd felépítése
\section{Érzékelés, érzékszervek}
% szem, látás & hallás, helyzetértékelés, szaglás, ízérzékelés
\section{A belső szervek szabályozása}
\section{Az ember egyedfejlődése} % maybe rename to Egyedfejlődés

% from here on its fun topics
% maybe these 3 should be 1 section and just subsections
\section{Kisfiú vagy kislány?} % maybe rename to Kromoszómák
\section{A tulajdonságok öröklődése}
\section{Nemhez kapcsolt öröklődés} 

% then random easy topics
% maybe these should also be 1 section "betegségek" and just subsection the rest
\section{Egészség, betegség}
\section{A diagnosztikai eljárások}
\section{A leggyakrabban előforduló betegségek}
\section{A szendvedélybetegségek}
\section{Balesetek, elsősegélynyújtás}

\end{document}